%% chapters/chapter04/chapter04.tex
%% CHAPTER 4: RESULTS AND ANALYSIS

\chapterfigpath{chapter04}

\chapter{Results and Analysis}
\label{chap:results}


\section{Overview of Results}
\label{sec:results_overview}

This chapter presents the results of the experiments conducted to evaluate the proposed solution. A total of 100 test cases were executed across multiple scenarios.


\section{Performance Evaluation}
\label{sec:performance_results}

Table~\ref{tab:performance_results} summarizes the performance measurements obtained during testing.

\begin{table}[H]
    \centering
    \caption{Performance measurement results}
    \label{tab:performance_results}
    \begin{tabular}{@{}lcccc@{}}
        \toprule
        \textbf{Test Case} & \textbf{Accuracy (\%)} & \textbf{Time (ms)} & \textbf{Score} & \textbf{Status} \\
        \midrule
        Test 1 & 95.2 & 120 & 4.5 & Pass \\
        Test 2 & 92.8 & 145 & 4.2 & Pass \\
        Test 3 & 88.5 & 180 & 3.8 & Pass \\
        Test 4 & 91.3 & 135 & 4.1 & Pass \\
        Test 5 & 94.7 & 125 & 4.4 & Pass \\
        \midrule
        \textbf{Average} & \textbf{92.5} & \textbf{141} & \textbf{4.2} & -- \\
        \bottomrule
    \end{tabular}
\end{table}

Figure~\ref{fig:results_comparison} provides a visual comparison of the results across different test scenarios.

\begin{figure}[ht]
    \centering
    \includegraphics[width=0.7\textwidth]{results_comparison}
    \caption{Comparison of results across test scenarios}
    \label{fig:results_comparison}
\end{figure}


\section{Functional Testing Results}
\label{sec:functional_results}

\begin{table}[H]
    \centering
    \caption{Functional testing summary}
    \label{tab:functional_tests}
    \begin{tabular}{@{}llcc@{}}
        \toprule
        \textbf{Module} & \textbf{Test Type} & \textbf{Total Tests} & \textbf{Pass Rate (\%)} \\
        \midrule
        User Authentication & Unit & 15 & 100.0 \\
        Data Processing & Unit & 28 & 96.4 \\
        API Endpoints & Integration & 20 & 95.0 \\
        User Interface & UI/UX & 12 & 91.7 \\
        \midrule
        \textbf{Overall} & -- & \textbf{75} & \textbf{95.8} \\
        \bottomrule
    \end{tabular}
\end{table}


\section{User Study Results}
\label{sec:user_study}

A user study was conducted with 30 participants. Table~\ref{tab:survey_results} shows their responses.

\begin{table}[H]
    \centering
    \caption{User survey results (Scale: 1-5, where 5 = Strongly Agree)}
    \label{tab:survey_results}
    \begin{tabular}{@{}lcccc@{}}
        \toprule
        \textbf{Question} & \textbf{Mean} & \textbf{Std Dev} & \textbf{Min} & \textbf{Max} \\
        \midrule
        The system is easy to use & 4.2 & 0.7 & 3 & 5 \\
        The interface is intuitive & 4.0 & 0.8 & 2 & 5 \\
        Results are accurate & 4.3 & 0.6 & 3 & 5 \\
        I would recommend this system & 4.1 & 0.9 & 2 & 5 \\
        \bottomrule
    \end{tabular}
\end{table}


\section{Comparison with Existing Solutions}
\label{sec:comparison}

Figure~\ref{fig:performance_graph} shows the performance comparison between the proposed solution and existing approaches.

\begin{figure}[ht]
    \centering
    \includegraphics[width=0.7\textwidth]{performance_graph}
    \caption{Performance comparison with existing solutions over time}
    \label{fig:performance_graph}
\end{figure}

Table~\ref{tab:comparison_results} provides a detailed comparison.

\begin{table}[H]
    \centering
    \caption{Comparison with existing solutions}
    \label{tab:comparison_results}
    \begin{tabular}{@{}lcccc@{}}
        \toprule
        \textbf{Feature} & \textbf{Proposed} & \textbf{System A} & \textbf{System B} & \textbf{System C} \\
        \midrule
        Accuracy (\%) & 92.5 & 85.0 & 88.2 & 90.1 \\
        Speed (ms) & 141 & 250 & 180 & 165 \\
        Cost & Free & \$99/mo & \$49/mo & Free \\
        Feature X & \checkmark & \checkmark & -- & \checkmark \\
        Feature Y & \checkmark & -- & \checkmark & -- \\
        \bottomrule
    \end{tabular}
\end{table}


\section{Statistical Analysis}
\label{sec:statistical_analysis}

A t-test was performed to compare the proposed system with existing solutions. The results showed a statistically significant difference ($t(28) = 2.45$, $p < 0.05$), indicating that the proposed approach outperforms the baseline.


\section{Discussion}
\label{sec:discussion}

\subsection{Key Findings}
\label{subsec:key_findings}

The experimental results reveal several important findings:

\begin{enumerate}
    \item \textbf{Finding 1:} The proposed system achieved 92.5\% accuracy, exceeding the target of 90\%
    
    \item \textbf{Finding 2:} Response time averaged 141ms, meeting the real-time requirement
    
    \item \textbf{Finding 3:} User satisfaction scores indicate high acceptance of the system
\end{enumerate}

\subsection{Addressing Research Questions}
\label{subsec:addressing_rq}

\textbf{RQ1:} The first research question was addressed through the performance evaluation, which demonstrated the effectiveness of the proposed approach.

\textbf{RQ2:} The second research question was answered by the comparison study, showing improvements over existing solutions.

\textbf{RQ3:} The third research question was addressed through the user study, which validated the usability of the system.

\subsection{Limitations of Results}
\label{subsec:results_limitations}

The results should be interpreted considering the following limitations:

\begin{itemize}
    \item Limitation 1 - e.g., ``Sample size was limited to 30 participants''
    \item Limitation 2 - e.g., ``Testing was performed only on synthetic data''
    \item Limitation 3 - e.g., ``Results may not generalize to other domains''
\end{itemize}


\section{Chapter Summary}
\label{sec:results_summary}

This chapter presented the results of the evaluation. Key findings include high accuracy (92.5\%), acceptable response time (141ms), and positive user feedback (4.2/5). The system achieved its main objectives while acknowledging certain limitations. These results address the research objectives outlined in Chapter 1.
